Este trabalho mostrou que estruturas de dados clássicas bastam para montar um recomendador completo para a tarefa de completar um look, e que cada etapa pode ser explicada. A KD-tree encontra peças parecidas sem comparar a consulta com todas as peças. O grafo de peças e looks, com BFS de dois passos, vai da peça aos looks dela e desses looks às outras peças, e assim traz o sinal dos looks reais. O Union-Find e o Kruskal formam os grupos de estilo. O algoritmo de Tarjan aponta as articulações entre esses grupos. As contribuições da equipe são quatro. A primeira são as versões A a F e a comparação entre elas. A segunda é uma avaliação justa, que pontua a categoria inteira, usa 10 ordens de desempate e informa intervalos de bootstrap. A terceira é a base tratada, publicada no Zenodo~\cite{zenodo}. A quarta é o código aberto~\cite{repositorio}. No teste, a versão final F chegou a Hit@10 de 1,056\% e FITB de 32,04\%, com ganho real sobre a versão C nas duas métricas. As medidas de tempo confirmaram na prática o que a teoria prevê para a KD-tree contra a força bruta e para o Union-Find com e sem a união por tamanho e a compressão de caminho.

O trabalho tem limitações. Em 72,6\% dos looks de teste, a peça escondida não aparece em nenhum look de treino. Nesse grupo, a F não coloca nenhuma escondida entre as 10 primeiras. Os atributos extraídos do texto são fracos. Na comparação com peças rotuladas pela imagem, a melhor fonte de cada atributo acerta de 51,3\% (estilo) a 77,0\% (estampa). O agrupamento pelo Kruskal sofre o efeito corrente, que concentra a maior parte das peças em poucos grupos. Além disso, os grupos não trouxeram ganho de acerto que se possa medir. A vantagem da F sobre a popularidade no Hit@10 também pode ser sorte.

Como próximas etapas, sugerimos extrair os atributos direto da imagem, para que a nota por atributos tenha força na peça nova. Sugerimos também identificar peças repetidas por um código de produto confiável, para ligar peças novas a peças já vistas. Outra direção é trocar o corte do Kruskal por um critério que evite o efeito corrente~\cite{zahn1971}. Por fim, a análise de articulações pode ser repetida em uma base com mais looks por peça.
